%!TEX root = ../main/main.tex

Regular expressions are one of the standard tools for processing textual data, and they play a central role in pattern matching, parsing, and information extraction. In the context of regex-based information extraction, the goal is not only to decide whether a pattern occurs in a document, but also to extract the relevant spans associated with variable captures. This perspective is naturally formalized through document spanners and related models such as regex formulas, ref-words, and variable-set automata.

Work on \emph{REmatch} showed that regex-based information extraction can be given an all-match semantics, allowing one to capture all relevant occurrences of a pattern rather than only those selected by the classical leftmost-longest behavior of standard regex engines. This viewpoint is particularly appealing in information extraction settings, where overlapping matches and nested captured spans often carry meaningful information.

However, for real-world text documents may contain typographical errors, missing symbols, extra insertions, or small local corruptions that prevent an exact match while preserving the intended structure of the information being queried. This raises a natural question: how should regex-based information extraction be extended to support \emph{errors}?

In this paper, we study two natural semantics for regex-based information extraction with errors. In the first one, matching is performed over a modified document obtained from the original input through a bounded sequence of edit operations. In the second one, matches are reported over the original document together with an edit script witnessing how the document can be transformed to satisfy the query. Although both viewpoints are intuitively natural, they differ in how spans are represented and in where the extracted mappings live: the edited document in the first case, and the original document in the second.

Our main goal is to show that these two views are formally equivalent. To do so, we first define how edit operations act not only on documents, but also on mappings, so that captured spans remain positionally consistent across transformations. We then study equivalence between edit scripts. This gives a clean representation of error witnesses and allows us to compare matches across the two semantics in a robust way.

Our main result is a bijective correspondence between the matches produced by both semantics. More precisely, we show that every match over a modified document can be translated into a match over the original document by transporting spans through the inverse edit sequence, and conversely, that every match over the original document induces a corresponding match over the edited document. This equivalence result provides a formal foundation for interoperating between systems that reason over edited inputs and systems that report extractions over the original document.
